\providecommand{\sacramentLexiconPageHook}{}
\newcommand{\lexrow}[2]{\textbf{#1} & #2\\}

\begin{landscape}
\hypersetup{pdfview=Fit}
\sacramentLexiconPageHook
\hbadness=10000
\thispagestyle{plain}
\section*{Metaphysical and Sacramental Lexicon}
{\small\noindent\emph{Classical philosophical and theological terms used analogically to distinguish divine agency, sensible sign, ecclesial mediation, and creaturely reception without reducing grace to a natural mechanism.}}

\vspace{0.25em}
\begingroup
\fontsize{7.15}{7.75}\selectfont
\setlength{\tabcolsep}{3pt}
\renewcommand{\arraystretch}{1.12}
\begin{tabularx}{\linewidth}{>{\RaggedRight\arraybackslash}p{0.205\linewidth}>{\RaggedRight\arraybackslash}X}
\toprule
\textbf{Term} & \textbf{Definition and sacramental bearing}\\
\midrule
\lexrow{Act, potency, and obediential potency}{Potency is real capacity for perfection; act is possessed perfection. \emph{Obediential potency} names a creature's radical openness to be elevated by God's omnipotence beyond every naturally due end, so grace perfects nature without becoming owed to it (\emph{ST}~I, q.~105, a.~6; III, q.~11, a.~1).}
\lexrow{Sacrament (\latin{sacramentum / mysterion})}{An efficacious sensible sign of invisible grace, instituted by Christ and entrusted to the Church, signifying Christ's Passion (commemorative), present sanctification (demonstrative), and future glory (prognostic), and causing the grace it signifies (\emph{ST}~III, q.~60, aa.~1--3; CCC~1131).}
\lexrow{Matter, quasi-matter, and form}{Hylomorphic analogy applied to the outward rite: \emph{matter} is the determinate sensible thing or gesture (water, chrism, wheat bread and grape wine, oil of the sick, imposition of hands); \emph{form} is the determinate word specifying its sanctifying meaning (\emph{ST}~III, q.~60, aa.~6--7). In Penance the penitent's acts serve as \emph{quasi-matter} (Trent, Sess.~XIV, ch.~3); in Matrimony mutual consent is assigned analogically.}
\lexrow{\latin{Sacramentum tantum}}{The outward sensible sign alone---matter and form enacted by a valid minister with ecclesial intention---which signifies and instrumentally effects the inward realities.}
\lexrow{\latin{Res et sacramentum}}{The intermediate, real spiritual effect produced immediately by a valid celebration even where an obstacle blocks grace: the indelible character in Baptism, Confirmation, and Holy Orders; Christ's true Body and Blood under the species in the Eucharist; the perpetual bond in Matrimony; inward penitential reconciliation with the Church in Penance; and bodily-spiritual consecration in Anointing. It is both reality (\latin{res}) of the outward rite and sign (\latin{sacramentum}) of the ultimate grace.}
\lexrow{\latin{Res tantum}}{The ultimate proper grace of each sacrament---sanctifying grace with its specific sacramental help (\emph{ST}~III, q.~62, a.~2), charity, ecclesial unity, healing, or spousal fidelity---which is reality signified and caused without itself serving as a further sacramental sign.}
\lexrow{Principal and instrumental causality}{God alone is the principal efficient cause of grace, working by His own power; Christ's sacred humanity is His \emph{conjoined instrument} (\latin{instrumentum coniunctum}); the sacramental rites and ordained ministers are \emph{separate instruments} (\latin{instrumenta separata}) through which the power of Christ's Passion reaches the soul (\emph{ST}~III, q.~62, aa.~1, 5).}
\lexrow{\latin{Ex opere operato} / \latin{ex opere operantis}}{The phrase \emph{ex opere operato} (``from the work worked'') teaches that a validly celebrated sacrament confers grace by Christ's own institution and power, independent of the minister's personal holiness, provided the recipient places no obstacle (Trent, Sess.~VII, can.~8). \emph{Ex opere operantis} names the fruit shaped by the dispositions, faith, and charity of the recipient or minister.}
\lexrow{Sacramental character}{An indelible spiritual power imprinted on the soul's intellective potency by Baptism, Confirmation, and Holy Orders, configuring the recipient to Christ the Priest and deputing him for divine worship (\emph{ST}~III, q.~63, aa.~1--6; Trent, Sess.~VII, can.~9). Character endures forever and forbids repetition of these three sacraments.}
\lexrow{Substance, accidents, and transubstantiation}{Substance is what exists in itself as a determinate kind of being; accidents (dimensions, color, taste, weight, nourishment) ordinarily inhere in a substance. In Eucharistic \emph{transubstantiation}, the entire substance of bread and wine is converted into the Body and Blood of Christ while the sensible appearances (\latin{species}) remain without their former substance, sustained miraculously by divine power (Lateran~IV, \emph{Firmiter}; Trent, Sess.~XIII, ch.~4, can.~2; \emph{ST}~III, qq.~75--77). It is neither a physical alteration of accidents nor identical with the words of consecration.}
\lexrow{Concomitance}{Because the risen Christ dies no more, wherever His Body is made present by the force of the sacramental words (\latin{ex vi verborum}), His Blood, soul, and divinity are wholly present by natural and hypostatic concomitance, and likewise under the species of wine (Trent, Sess.~XIII, ch.~3; \emph{ST}~III, q.~76, aa.~1--2; CCC~1390).}
\lexrow{Minister and intention}{Christ is the primary minister of every sacrament; human ministers act in His name as intentional instruments. Validity requires at least the \emph{intention to do what the Church does} (Trent, Sess.~VII, can.~11; \emph{ST}~III, q.~64, aa.~8--10). Law distinguishes ordinary from extraordinary ministers and, in Penance, faculties of jurisdiction alongside valid priestly Orders.}
\lexrow{Subject, \latin{obex}, and reviviscence}{The \emph{subject} is the capable human recipient (unbaptized for Baptism; validly baptized for the other six). An \emph{obex} is an interior obstacle---such as unrepented mortal sin or lack of faith---that blocks fruitful grace without invalidating the sacrament. Upon later sincere repentance, grace blocked at reception may take effect through \emph{reviviscence} in character-imprinting and enduring-bond sacraments.}
\lexrow{Validity, liceity, and fruitfulness}{\emph{Validity} concerns whether the sacrament truly exists (essential sign, valid minister, ecclesial intention, capable subject); \emph{liceity} concerns lawful celebration according to Church discipline; \emph{fruitfulness} concerns the actual reception and increase of grace in a well-disposed recipient.}
\lexrow{Order of ends}{Classical analysis distinguishes a sacrament's \emph{primary proper end} (its defining intrinsic purpose), \emph{intrinsic secondary ends} (real goods belonging to the institution itself), \emph{contingent ends} (circumstantial or pastoral benefits), and \emph{ultimate end} (the glory of the Triune God and eternal beatitude).}
\lexrow{Church \latin{sui iuris} and rite}{A \emph{Church sui iuris} is a community of Christian faithful linked by hierarchy and recognized by supreme ecclesiastical authority (CCEO, c.~27); twenty-four stand in Catholic communion (the Latin Church and twenty-three Eastern Catholic Churches). A \emph{rite} is the liturgical, theological, spiritual, and disciplinary patrimony proper to each Church across the Latin and five Eastern traditions (CCEO, c.~28).}
\lexrow{Sacramentals}{Sacred signs instituted by the Church (blessings, consecrations, exorcisms, liturgical gestures) that bear a likeness to the sacraments, signify chiefly spiritual effects obtained through the Church's intercession (\latin{ex opere operantis Ecclesiae}), and dispose the faithful to receive the chief effect of the sacraments (\emph{SC}~60).}
\bottomrule
\end{tabularx}
\endgroup
\end{landscape}
